\section{Error Measures}

\subsection{Binary Classification}
A binary classification model predicts the probability \f{p(y_n = 1 | x_n) \approx f(x_n; \theta)}. The final estimated target \f{\hat{y}_n} is determined by applying a threshold \f{\epsilon} to the predicted probability. The results yield a confusion matrix.
%  consisting of True Positives (TP), False Positives (FP), True Negatives (TN), and False Negatives (FN).


\definition{Misclassification Rate (MCR)}{The ratio of incorrect predictions to the total number of predictions, which can be expressed mathematically as a sum of 0/1 error terms (0/1 step function is not differentiable, so does not work with SGD).
\cf{MCR(y, \hat{y}) = \frac{FP + FN}{TP + TN + FP + FN} = \frac{1}{N}\sum_{n=1}^N \mathbb{I}[y_n \neq \hat{y}_n]}
\b{Important}: MCR is a highly misleading metric for imbalanced datasets (e.g., rare disease detection) because predicting only the majority class yields a high accuracy while completely failing the actual task.
}

\definition{Proxy Loss}{A differentiable function that approximates the true, non-differentiable error metric. Minimizing the proxy loss should roughly lead to the same optimal parameters as minimizing the true metric.\\[.5em]
\b{Example}: BCE acts as a smooth proxy for 0/1 MCR error.}



\subsection{Advanced Classification Metrics}

\begin{itemize}
\item \b{Recall (True Positive Rate)}: Ratio of correctly detected positive cases out of all actual positive cases.
\cf{TPR = \frac{TP}{TP + FN}}
\item \b{False Positive Rate (FPR)}: Ratio of false alarms out of all actual negative cases.
\cf{FPR = \frac{FP}{TN + FP}}
\item \b{Receiver Operating Characteristic (ROC)}: A curve that plots TPR against FPR across all possible decision thresholds \f{\epsilon \in [0, 1]}.
\item \b{AUC}: The Area Under the ROC Curve, providing an aggregate measure of performance across all thresholds. It is maximized by correctly ranking positive instances higher than negative ones.
\item \b{Precision}: The ratio of actual positive cases out of all cases that the model predicted as positive.
\cf{Precision = \frac{TP}{TP + FP}}
\item \b{F1 Score}: The harmonic mean of Precision and Recall. It is heavily used for imbalanced datasets. Because F1 is non-differentiable and non-decomposable, a common workaround is to oversample the minority class, train the model using a proxy loss, and then tune hyperparameters to explicitly maximize the validation F1 score.
\cf{F1 = 2 \frac{Precision \cdot Recall}{Precision + Recall}}
\end{itemize}
\vspace{1em}
\cig{.25}{roc}




\subsection{Regression Error Measures}
Regression error measures are naturally differentiable and do not require proxies.
\begin{itemize}
\item \b{Scale-dependent errors}: Include Mean Absolute Error (MAE) and Root Mean Square Error (RMSE).
\end{itemize}

\cf{RMSE = \sqrt{\frac{1}{N}\sum_{n=1}^N (y_n - \hat{y}_n)^2}}

\begin{itemize}
\item \b{Scale-independent (percentage) errors}: Include Mean Absolute Percentage Error (MAPE). However, MAPE becomes undefined or extreme when the true target \f{y_n \approx 0}.
\end{itemize}

\cf{MAPE = \frac{1}{N}\sum_{n=1}^N \left|\frac{y_n - \hat{y}_n}{y_n}\right|}

\begin{itemize}
\item \b{Mean Absolute Scaled Error (MASE)}: Fixes the instability of MAPE by scaling the error using the average error of a naive baseline model instead of the true target value.
\end{itemize}











\subsection{Ranking \& Information Retrieval}
The goal in ranking is to learn a function that outputs relevance scores \f{s_i = f(x_i)} for documents matching a specific query \f{q}, such that the predicted order matches the true relevance labels \f{l_i}.
\begin{itemize}
\item \b{Discounted Cumulative Gain (DCG@K)}: A measure that sums the true relevances of the top K ranked documents, heavily penalizing relevant documents that are placed lower in the ranking.
\cf{DCG@K = \sum_{i=1}^K \frac{2^{l_i} - 1}{\log_2(i + 1)}}
\item \b{Normalized DCG (NDCG@K)}: DCG scales unboundedly depending on the query. NDCG normalizes it to a \f{[0, 1]} range by dividing the predicted DCG by the Ideal DCG (IDCG), which is the DCG obtained by a perfectly sorted ranking.
\cf{NDCG@K = \frac{DCG@K}{IDCG@K}}
\end{itemize}


\definition{Pairwise Rank Loss}{A technique that reframes the ranking problem by estimating the probability that any given pair of documents \f{(i, j)} is correctly ranked relative to each other.\\[1em]
To train a pairwise ranking model, we define a target state \f{S_{i,j}} (\f{1} if document \f{i} is more relevant, \f{-1} if less, and \f{0} if equal) and optimize the parameters to minimize the pairwise log loss.
\cf{\mathcal{L}*{i,j} = \frac{1}{2}(1 - S*{i,j})(s_i - s_j) + \log(1 + e^{-(s_i - s_j)})}
}

\vfill
\cig{.98}{methods}